% Relatorio de Projeto IF674 2026.2 — equipe EQ01
% Compilacao sugerida: pdflatex relatorio.tex (duas vezes para atualizar indices).
% Edite os dados da capa; as imagens do Figma estao em figuras/.
\documentclass[12pt,a4paper]{article}

\usepackage[T1]{fontenc}
\usepackage[utf8]{inputenc}
\usepackage[brazilian]{babel}
\usepackage{lmodern}
\usepackage{geometry}
\geometry{left=0.8cm,right=0.8cm,top=2.6cm,bottom=1cm,footskip=0.55cm}
\usepackage{microtype}
\usepackage{graphicx}
\usepackage{pdflscape}
\usepackage{float}
\usepackage{array}
\usepackage{tabularx}
\usepackage{longtable}
\usepackage{booktabs}
\usepackage{xcolor}
\usepackage{tikz}
\usetikzlibrary{arrows.meta,positioning,shapes.geometric,fit}
\usepackage{amsmath}
\usepackage{fancyhdr}
\usepackage{hyperref}
\hypersetup{hidelinks,pdftitle={Relatorio de Projeto IF674 2026.2},pdfauthor={Carlos Henrique Goncalves dos Santos; Robson Thiago Duarte}}
\graphicspath{{figuras/}}
\setlength{\parindent}{1.25cm}
\setlength{\parskip}{0.25em}
\renewcommand{\arraystretch}{1.2}
\emergencystretch=2em

\pagestyle{fancy}
\fancyhf{}
\fancyhead[L]{IF674 — Relatório de Projeto}
\fancyhead[R]{Equipe EQ01}
\fancyfoot[C]{\thepage}
\renewcommand{\headrulewidth}{0.3pt}
\setlength{\headheight}{15pt}

% Preencha apenas a data oficial da entrega antes de submeter.
\newcommand{\Equipe}{EQ01}
\newcommand{\Integrantes}{Carlos Henrique Gonçalves dos Santos (CC) --- \texttt{chgs3}\\[0.25cm]Robson Thiago Duarte (EC) --- \texttt{rtds}}
\newcommand{\DataEntrega}{\textit{Preencher a data oficial da entrega}}

% A imagem aparece automaticamente ao ser salva com o nome indicado em figuras/.
% Enquanto ela nao existe, o PDF mostra um campo reservado com a mesma legenda.
\newcommand{\FiguraFigma}[4]{%
  \begin{figure}[H]
  \centering
  \IfFileExists{figuras/#1}{%
    \includegraphics[width=#4\linewidth,keepaspectratio]{figuras/#1}%
  }{%
    \fbox{\begin{minipage}[c][5.2cm][c]{0.9\linewidth}
      \centering\textbf{Imagem a inserir}\par\medskip
      Exporte do Figma e salve como \texttt{figuras/#1}.\par\medskip
      #2
    \end{minipage}}%
  }%
  \caption{#2}\label{#3}
  \end{figure}
}

% Modelo repetivel. O conteudo entre as linhas e propositalmente aberto.
\newcommand{\SimulacaoInstrucao}[4]{%
  \subsection{#1}\label{sim:#2}
  \noindent\textbf{Objetivo:} #3\par
  \noindent\textbf{Preparação e programa de teste:} \textit{Inserir valores iniciais, palavra de instrução e sequência usada.}\par
  \noindent\textbf{Sinais observados:} clock, reset, estado da UC, PC, EPC, IR, MDR, R29, R31 e os registradores e sinais específicos da instrução.\par
  \noindent\textbf{Resultado esperado:} #4\par
  \noindent\textbf{Resultado observado e análise ciclo a ciclo:} \textit{Inserir os dados da simulação realizada. Descrever estados, valores antes e depois das bordas e comparação com o resultado esperado.}\par
  \begin{figure}[H]\centering
    \fbox{\begin{minipage}[c][3.5cm][c]{0.9\linewidth}\centering
      Inserir captura legível da forma de onda Quartus para #1.
    \end{minipage}}
    \caption{Simulação de #1 — inserir captura após validação.}\label{fig:sim-#2}
  \end{figure}
}

\begin{document}
\hypersetup{pageanchor=false}
\begin{titlepage}
  \centering
  {\large UNIVERSIDADE FEDERAL DE PERNAMBUCO\par}
  {\large CENTRO DE INFORMÁTICA\par}
  \vspace{1.1cm}
  {\large INFRA-ESTRUTURA DE HARDWARE — IF674\par}
  \vfill
  {\LARGE\bfseries Relatório de Projeto\par}
  \vspace{0.6cm}
  {\large Processador multiciclo de 32 bits — 2026.2\par}
  \vfill
  {\large Equipe \Equipe\par}
  \vspace{0.5cm}
  \Integrantes\par
  \vfill
  Recife — 21/10/2026\par
\end{titlepage}

\hypersetup{pageanchor=true}
\pagenumbering{roman}
\tableofcontents
\clearpage
\listoffigures
\clearpage
\listoftables
\clearpage
\pagenumbering{arabic}

\section{Introdução}\label{sec:introducao}

Este relatório descreve a arquitetura e o controle de um processador multiciclo de 32 bits para as 23 instruções definidas na especificação do projeto de IF674 de 2026.2~\cite{especificacao}. O objetivo é associar o comportamento de cada instrução aos blocos do caminho de dados, aos sinais da unidade de controle (UC) e à sequência temporal de estados. A versão 2026.2 substitui \texttt{xchg} e \texttt{sram} por \texttt{push} e \texttt{pop}; por isso, o tratamento da pilha e do registrador R29 recebe atenção especial nas Seções~\ref{sec:pilha} e~\ref{sec:estados-pilha}.

Por decisão de Robson Thiago Duarte, as instruções são apresentadas em agrupamentos definidos pela semelhança de seus caminhos de dados e etapas de controle. Essa organização evidencia os recursos compartilhados da arquitetura multiciclo; o catálogo individual das 23 instruções permanece no Apêndice~\ref{ap:catalogo}, e o roteiro de simulação de cada uma está na Seção~\ref{sec:simulacoes}.

O projeto foi organizado a partir da documentação das partes 1, 2 e 3 da equipe~\cite{docpt1,docpt2,docpt3}, do quadro no FigJam~\cite{figma} e dos conteúdos de instruções, processadores monociclo e multiciclo e desempenho trabalhados na disciplina~\cite{aulas}. A arquitetura multiciclo reutiliza unidades em ciclos diferentes e preserva valores intermediários em IR, A, B, ALUOut e MDR. A UC determina quais fontes, operações e escritas ficam ativas em cada estado. Essa distinção entre caminho de dados e controle orienta tanto os fluxos deste texto como os espaços preparados para as formas de onda.

\subsection{Escopo e critérios do projeto}\label{subsec:escopo}

A CPU possui 32 registradores gerais de 32 bits e uma memória unificada de 256 bytes endereçada por byte. O conjunto tem 13 instruções R, oito I e duas J, catalogadas no Apêndice~\ref{ap:catalogo}. A memória lê de forma síncrona, escreve em um ciclo e armazena bytes em ordem \emph{little-endian}. O registrador R29 atua como ponteiro da pilha (SP), com valor inicial especificado em 227. A UC também deve tratar opcode inexistente, overflow e divisão por zero, preservando em EPC o endereço da instrução causadora~\cite{especificacao}.

O Capítulo 3 da especificação exige capa, índice, introdução, unidade de processamento, descrição das entidades, descrição dos estados de controle, conjunto de simulações e conclusão. A estrutura deste relatório segue essa ordem. As exportações do circuito e do mapa de estados estão nas Figuras~\ref{fig:circuito-figma} e~\ref{fig:mapa-figma}; os fluxos menores são desenhados diretamente em LaTeX para permanecerem nítidos e fáceis de atualizar.

\subsection{Situação da implementação e da evidência}\label{subsec:situacao}

Os documentos disponíveis registram o desenho, as decisões de caminho de dados e os estados propostos de \texttt{push}/\texttt{pop}. A parte 3 acrescenta seis módulos VHDL fornecidos, o arquivo inicial de memória e onze programas Assembly de apoio~\cite{docpt3}. Esses artefatos permitem conferir interfaces e preparar estímulos; a integração da CPU e da UC e as capturas Quartus das 23 instruções ainda precisam ser documentadas neste relatório. A Seção~\ref{sec:simulacoes} contém um roteiro individual para cada instrução e para as entidades implementadas; os campos de resultado observado permanecem abertos para receber as evidências. As pendências técnicas de integração e orientação da disciplina aparecem no Apêndice~\ref{ap:pendencias}.

\subsection{Relação com as aulas}\label{subsec:aulas}

As aulas sobre representação de instruções fundamentam a decodificação dos campos \texttt{opcode}, \texttt{rs}, \texttt{rt}, \texttt{rd}, \texttt{shamt}, \texttt{funct} e imediato. A implementação monociclo ajuda a identificar o efeito completo de uma instrução; a implementação multiciclo reparte esse efeito em etapas que compartilham ALU, memória e banco de registradores. A aula de desempenho mostra que a economia de hardware deve ser analisada junto com número de ciclos, período do clock e mistura de instruções. Portanto, a contagem de caixas em um fluxo não basta, sozinha, para concluir qual arquitetura é mais rápida~\cite{aulas}.

\section{Unidade de processamento}\label{sec:processamento}

\subsection{Visão geral do caminho de dados}\label{subsec:datapath}

O caminho de dados reúne três funções: buscar a instrução e atualizar PC; ler operandos e executar a operação; acessar a memória e escrever o resultado. A memória é compartilhada entre busca e dados. IorD escolhe PC durante a busca e ALUOut durante cargas, armazenamentos e operações de pilha. IR conserva a instrução; A e B conservam as leituras do banco; ALUOut retém endereço efetivo ou resultado; MDR retém a saída de leitura da memória.

\begin{figure}[H]
\centering
\begin{tikzpicture}[>=Latex, node distance=4mm,
  bloco/.style={draw,rounded corners,align=center,minimum height=10mm,text width=18mm,font=\footnotesize,inner sep=2pt},
  controle/.style={draw,rounded corners,align=center,minimum height=8mm,text width=65mm,font=\footnotesize,fill=gray!10}]
  \node[bloco] (pc) {PC};
  \node[bloco,right=of pc] (iord) {IorD};
  \node[bloco,right=of iord] (mem) {memória};
  \node[bloco,right=of mem] (ir) {IR};
  \node[bloco,right=of ir] (banco) {banco};
  \node[bloco,right=of banco] (ab) {A e B};
  \node[bloco,right=of ab] (alu) {ALU};
  \node[bloco,right=of alu] (aluout) {ALUOut};
  \node[controle,above=10mm of ir] (uc) {UC: seletores, habilitações e próximo estado};
  \foreach \a/\b in {pc/iord,iord/mem,mem/ir,ir/banco,banco/ab,ab/alu,alu/aluout}{\draw[->] (\a) -- (\b);}
  \draw[->,dashed] (uc.south) -- (ir.north);

  \node[bloco,below=16mm of iord] (endereco) {ALUOut};
  \node[bloco,right=of endereco] (iorddados) {IorD};
  \node[bloco,right=of iorddados] (memdados) {memória};
  \node[bloco,right=of memdados] (mdr) {MDR};
  \node[bloco,right=of mdr] (wb) {MemtoReg};
  \node[bloco,right=of wb] (bancodados) {banco};
  \foreach \a/\b in {endereco/iorddados,iorddados/memdados,memdados/mdr,mdr/wb,wb/bancodados}{\draw[->] (\a) -- (\b);}
  \node[font=\scriptsize,below=1mm of pc] {busca};
  \node[font=\scriptsize,below=1mm of endereco] {dados};
\end{tikzpicture}
\caption{Fluxos funcionais da CPU multiciclo. Cada linha resume uma sequência de estados; a UC configura os caminhos a cada ciclo.}\label{fig:fluxo-geral}
\end{figure}

A Figura~\ref{fig:fluxo-geral} serve para leitura conceitual. O desenho completo da equipe, com seletores e sinais da UC, aparece na Figura~\ref{fig:circuito-figma}. O recorte seguinte amplia as conexões centrais do caminho de dados.

\FiguraFigma{circuito-completo.png}{Circuito completo do processador multiciclo exportado do FigJam.}{fig:circuito-figma}{0.98}

O recorte da Figura~\ref{fig:circuito-detalhe} amplia a região de busca, leitura de operandos, ALU e seleção de escrita. Ele complementa a visão integral da Figura~\ref{fig:circuito-figma}, sem substituir as demais ligações do quadro.

\FiguraFigma{circuito-detalhe-busca-alu.png}{Detalhe do circuito: memória, IR, banco, ALU, EPC e seletores próximos.}{fig:circuito-detalhe}{0.95}

\subsection{Formatos, busca e acesso à memória}\label{subsec:busca}

Os formatos R, I e J contêm, respectivamente, os campos \texttt{opcode|rs|rt|rd|shamt|funct},\linebreak \texttt{opcode|rs|rt|imediato} e \texttt{opcode|target}. A busca apresenta o endereço inicial $P$ à memória e conserva a instrução em IR quando a saída estiver válida. Em paralelo ao caminho de leitura, a ALU produz $P+4$ para o fluxo sequencial. O acesso posterior a dados usa o endereço guardado em ALUOut, sem confundir esse ciclo com a busca.

Para \texttt{lw}, \texttt{lb}, \texttt{sw} e \texttt{sb}, o endereço efetivo é
\begin{equation}\label{eq:endereco-efetivo}
  EA = R[rs] + \operatorname{SE}_{16}(\textit{offset}).
\end{equation}
Como a leitura da memória só fica disponível após a requisição, a UC deve manter endereço e IorD estáveis até o dado poder ser capturado em IR ou MDR. A escrita ocupa um ciclo. O módulo \texttt{Memoria.vhd} fornecido usa quatro bancos de oito bits e somente \texttt{Address(7 downto 0)} para os 256 bytes; a seleção e a rotação dos bancos acomodam os deslocamentos do endereço. Na ordem \emph{little-endian}, a palavra \texttt{0xAABBCCDD} iniciada no endereço decimal 223 corresponde aos bytes \texttt{DD}, \texttt{CC}, \texttt{BB} e \texttt{AA} nos endereços 223 a 226. Assim, a estrutura do componente contempla endereço de palavra desalinhado, embora o resultado e a latência precisem ser confirmados em simulação~\cite{especificacao,docpt3}.

\subsection{Grupos de reaproveitamento}\label{subsec:grupos}

\begin{table}[H]\centering\small
\caption{Famílias de instruções organizadas pelo caminho predominante.}\label{tab:grupos}
\begin{tabularx}{\linewidth}{@{}p{1.2cm}p{5.3cm}X@{}}\toprule
Grupo & Instruções & Caminho principal \\\midrule
G0 & Todas & PC, memória, IR e decodificação. \\
G1 & add, and, sub, slt, addi & Banco, ALU/comparador e escrita no banco. \\
G2 & sll, sra, lui & Fonte, SHIFT e escrita no banco. \\
G3 & lw, lb & Endereço, memória, MDR e escrita no banco. \\
G4 & sw, sb & Endereço, dado do banco e escrita na memória. \\
G5 & beq, bne & Comparação e escolha condicional do PC. \\
G6 & j, jal, jr & Formação/seleção do alvo e escrita do PC. \\
G7 & mult, div, mfhi, mflo & Unidades especiais e registradores HI/LO. \\
G8 & push, pop & R29, ALUOut, memória, MDR e escrita no banco. \\\bottomrule
\end{tabularx}
\end{table}

G0 é transversal: não adiciona instruções à contagem. Os grupos G1 a G8 particionam as 23 operações para explicar o reaproveitamento, mas não substituem a decodificação por \texttt{opcode} e \texttt{funct}. Nas aulas de multiciclo, a presença de registradores intermediários permite que o mesmo recurso físico atenda grupos diferentes em ciclos sucessivos~\cite{aulas}.

\subsection{Pilha e extensão da arquitetura}\label{sec:pilha}

O SP é R29. A pilha cresce para endereços menores. O conteúdo original do registrador fonte deve ser conservado antes de atualizar R29:
\begin{align}
\text{\texttt{push rs:}}\quad & d \gets R[rs]_{\mathrm{antigo}}, & SP &\gets SP_{\mathrm{antigo}}-4, & M_{32}[SP] &\gets d; \label{eq:push}\\
\text{\texttt{pop rd:}}\quad & d \gets M_{32}[SP_{\mathrm{antigo}}], & R[rd] &\gets d, & SP &\gets SP_{\mathrm{antigo}}+4. \label{eq:pop}
\end{align}

Na captura de \texttt{push}, Read1Sel seleciona R29 para A, e Read2Sel seleciona \texttt{rs} para B. Em \texttt{pop}, A recebe R29 e B não é necessário. RegDst ganha a opção R29; MemtoReg ganha a constante 227 para a proposta de inicialização; DataSource usa B como dado de escrita da memória. A largura de MemtoReg deve comportar a entrada de índice 8. Os sinais AWrite e BWrite são desabilitados depois da captura, para preservar o SP e o dado antigos durante toda a operação.

\begin{table}[H]\centering\small
\caption{Seleções essenciais da extensão da pilha.}\label{tab:sinais-pilha}
\begin{tabularx}{\linewidth}{@{}p{3cm}p{4.3cm}X@{}}\toprule
Sinal & Seleção relevante & Finalidade \\\midrule
Read1Sel & 1 = R29 & Capturar SP antigo em A. \\
Read2Sel & 1 = rs & Capturar dado de \texttt{push} em B. \\
RegDst & 1 = rd; 3 = R29 & Escrever destino de \texttt{pop} ou SP. \\
MemtoReg & 0 = ALUOut; 1 = MDR; 8 = 227 & Escrever endereço calculado, dado lido ou SP inicial. \\
DataSource & 2 = B & Escrever o dado antigo de \texttt{push}; também corrige a origem de \texttt{sw}. \\
IorD & 1 = ALUOut & Apresentar endereço de dados à memória. \\
ALUSrcA / ALUSrcB & A / constante 4 & Calcular $SP-4$ ou $SP+4$. \\\bottomrule
\end{tabularx}
\end{table}

\subsection{Demais operações e cuidados de integração}\label{subsec:demais}

\texttt{add}, \texttt{and}, \texttt{sub}, \texttt{slt} e \texttt{addi} compartilham leitura, execução e escrita no banco. \texttt{slt} deve produzir a palavra 0 ou 1, com comparação com sinal. \texttt{sll}, \texttt{sra} e \texttt{lui} compartilham SHIFT, mas variam fonte, quantidade e modo. Em \texttt{RegDesloc.vhd}, carregar a entrada e deslocá-la são comandos síncronos distintos; a UC deve prever ciclos separados ou demonstrar um caminho alternativo que produza o mesmo efeito. \texttt{lb} precisa selecionar um byte e estendê-lo com sinal; \texttt{sb} precisa preservar os três bytes vizinhos porque a memória não tem habilitação individual de byte.

\texttt{beq} e \texttt{bne} usam o alvo $P+4+(\operatorname{SE}_{16}(\textit{offset})\ll 2)$ e condições opostas. \texttt{j} e \texttt{jal} formam o alvo com os quatro bits altos do PC indicado pela especificação e \texttt{target}\,$\ll 2$; \texttt{jr} usa diretamente R[rs]. A documentação da equipe assume provisoriamente link $P+4$ e ausência de \emph{delay slot} em \texttt{jal}; essa convenção exige confirmação. MULT e DIV devem ser unidades separadas da ALU. A multiplicação sinalizada usa Booth, produz 64 bits e aceita conclusão entre 32 e 40 ciclos; HI/LO guardam as metades. Na divisão, HI guarda o resto, LO o quociente e a divisão por zero comunica exceção à UC~\cite{especificacao}.

\section{Descrição das entidades}\label{sec:entidades}

Esta seção especifica a função e as interfaces relevantes dos blocos. A descrição final de cada entidade construída pela equipe deverá refletir as portas e o algoritmo do HDL efetivamente entregue, como exige a Seção 3.2.5 da especificação~\cite{especificacao}. Componentes fornecidos pela disciplina são descritos aqui para explicar sua integração; sua simulação isolada não é exigida pelo capítulo do relatório.

\subsection{Componentes recebidos na parte 3}\label{subsec:componentes-pt3}

A pasta \texttt{Componentes do Projeto - 2026.2} reúne os seis módulos da Tabela~\ref{tab:modulos-pt3}. O arquivo \texttt{LEIA-ME.txt} orienta incluí-los no projeto sem alterar seu código e posicionar \texttt{instrucoes.mif} junto ao arquivo \texttt{.mpf}; o conteúdo do MIF pode ser atualizado para os programas e dados dos ensaios~\cite{docpt3}. Esses arquivos são componentes de base, e não constituem por si uma CPU integrada.

\begin{table}[H]\centering\small
\caption{Módulos VHDL e arquivo de memória fornecidos na parte 3.}\label{tab:modulos-pt3}
\begin{tabularx}{\linewidth}{@{}p{3.4cm}X X@{}}\toprule
Arquivo & Função observada no código & Ponto de integração \\\midrule
\texttt{Banco\_reg.vhd} & 32 registradores; duas leituras combinacionais e escrita síncrona. & Reset assíncrono zera também R29; conferir seleção e escrita de SP. \\
\texttt{ula32.vhd} & ALU combinacional de 32 bits com seletor de três bits. & Compatibilizar operação, flags e código da UC. \\
\texttt{RegDesloc.vhd} & Registro, deslocamentos lógicos e aritmético e rotações. & Carga e deslocamento usam bordas distintas. \\
\texttt{Registrador.vhd} & Registrador de 32 bits com habilitação de carga. & Usar para retenção temporal dos valores necessários. \\
\texttt{Instr\_Reg.vhd} & Registra instrução e expõe campos para decodificação. & Conferir extração posterior de \texttt{rd} e \texttt{funct}. \\
\texttt{Memoria.vhd} & Memória em quatro bancos de bytes, com leitura em clock. & Verificar ordem dos bytes, rotação e latência dos dados. \\
\texttt{instrucoes.mif} & Inicialização de 256 endereços de oito bits. & Carregar programa e dados correspondentes a cada ensaio. \\\bottomrule
\end{tabularx}
\end{table}

\subsection{Registradores e banco}\label{subsec:registradores}

PC, IR, A/B, ALUOut, MDR e EPC conservam, respectivamente, o endereço de busca, a instrução, os operandos, o resultado intermediário, o dado lido e o endereço da instrução que causou exceção. As habilitações de escrita permitem reutilizar os mesmos blocos em ciclos diferentes.

O banco possui 32 registradores de 32 bits, duas leituras combinacionais e uma escrita na borda de subida. O reset assíncrono de \texttt{Banco\_reg.vhd} zera também R29, o que exige conciliação com o SP inicial de 227 (Apêndice~\ref{ap:pendencias}). R31 recebe o retorno de \texttt{jal} conforme a convenção final adotada.

\subsection{ALU, deslocador e unidades especiais}\label{subsec:alu}

A ALU combinacional executa passagem de A (000), soma (001), subtração (010), AND (011), incremento (100), NOT (101) e XOR (110), além de fornecer flags de condição. A UC deve escolher o código em cada estado; \texttt{pop}, por exemplo, usa operações distintas para obter o endereço e incrementar SP. Em \texttt{RegDesloc.vhd}, carga (001), deslocamento à esquerda (010) e deslocamento aritmético à direita (100) ocorrem em bordas distintas, fato a considerar em \texttt{sll}, \texttt{sra} e \texttt{lui}.

As unidades MULT e DIV ainda dependem da implementação integrada. MULT deve produzir o produto de 64 bits por Booth com sinal; DIV deve produzir quociente e resto e sinalizar divisão por zero. HI/LO guardam os resultados, e a UC deve aguardar a conclusão antes de liberar sua leitura.

\subsection{Memória e tratamento de bytes}\label{subsec:memoria}

\texttt{Memoria.vhd} usa os oito bits menos significativos de \texttt{Address} para acessar 256 bytes em quatro bancos. A rotação dos bytes e seus registros de atraso exigem conferir em forma de onda a borda correta de captura em IR e MDR, inclusive para endereços desalinhados~\cite{docpt3}.

\texttt{lb} seleciona e estende com sinal o byte endereçado. \texttt{sb} requer leitura, mescla e reescrita para preservar os outros três bytes, pois a memória não possui \emph{byte enable}; esse caminho ainda precisa ser demonstrado no HDL e nas formas de onda.

\subsection{Unidade de controle e seletores}\label{subsec:uc-entidade}

A UC combina os campos de IR com flags, exceções e sinais de conclusão para escolher o próximo estado. Seus seletores definem as fontes do caminho de dados; suas habilitações controlam as escritas. O diagrama final deve identificar os sinais ativos em cada estado e as opções de cada multiplexador.


\section{Descrição dos estados do controle}\label{sec:controle}

\subsection{Como ler o mapa de estados}\label{subsec:leitura-mapa}

Cada estado representa uma etapa temporal da UC. Dentro de um estado, seletores e unidades combinacionais podem atuar simultaneamente; na borda, apenas os elementos habilitados capturam valores. As setas representam a sequência entre ciclos. Para explicar uma caixa do mapa, deve-se indicar: valores disponíveis no início, sinais ativados, cálculo durante o ciclo, efeito na borda e próximo estado. Essa regra, desenvolvida no guia da parte 2~\cite{docpt2}, impede confundir ações simultâneas com estados sucessivos.

O mapa completo da Figura~\ref{fig:mapa-figma} aparece ao fim desta seção em página paisagem; a Figura~\ref{fig:mapa-recorte} amplia a sequência dos estados. A legenda da Figura~\ref{fig:legenda-mapa} reúne os agrupamentos de sinais. O esquema resumido da Figura~\ref{fig:fluxo-estados} serve à navegação. O nome e a codificação binária dos estados desenhados devem ser conciliados com a UC implementada; aqui são usados nomes funcionais para não atribuir ao HDL estados ainda não conferidos.

\begin{figure}[H]\centering
\begin{tikzpicture}[>=Latex,node distance=6mm and 8mm,
  st/.style={draw,rounded corners,align=center,minimum height=8mm,text width=18mm,font=\scriptsize},
  stack/.style={draw,rounded corners,align=center,minimum height=8mm,text width=31mm,font=\scriptsize,fill=gray!10}]
\node[st] (init) {INIT\_SP};
\node[st,right=of init] (busca) {BUSCA};
\node[st,right=of busca] (dec) {DECODE};
\node[stack,above right=7mm and 13mm of dec] (familias) {ALU, SHIFT, memória,\\ desvios, HI/LO};
\node[stack,right=13mm of dec] (push) {PUSH\_DEC $\to$ PUSH\_SP\\ $\to$ PUSH\_STORE};
\node[stack,below right=7mm and 13mm of dec] (pop) {POP\_ADDR $\to$ REQ $\to$ WAIT\\ $\to$ CAP $\to$ RD $\to$ INC $\to$ SP};
\draw[->] (init) -- (busca);
\draw[->] (busca) -- (dec);
\draw[->] (dec) -- (familias);
\draw[->] (dec) -- (push);
\draw[->] (dec) -- (pop);
\node[font=\scriptsize,below=21mm of dec] {Cada ramo retorna à BUSCA ao terminar.};
\end{tikzpicture}
\caption{Resumo lógico da transição entre busca, decodificação e famílias de execução.}\label{fig:fluxo-estados}
\end{figure}

\FiguraFigma{legenda-mapa-estados.png}{Legenda de grupos e sinais de controle usada no mapa de estados da equipe.}{fig:legenda-mapa}{0.86}

\subsection{Estados compartilhados}\label{subsec:estados-comuns}

BUSCA apresenta PC à memória, aguarda a saída síncrona quando necessário, registra IR e calcula o PC sequencial. DECODE examina \texttt{opcode} e, para instruções R, \texttt{funct}; captura operandos e escolhe a família. Estados de execução da ALU registram resultados em ALUOut e os estados de escrita escolhem destino e dado. As cargas calculam EA pela Equação~\ref{eq:endereco-efetivo}, aguardam a leitura em MDR e escrevem no banco. Os armazenamentos apresentam EA e o dado adequado à memória. Os desvios atualizam PC somente quando sua condição se confirma; saltos escolhem o alvo formado ou R[rs]. Todos retornam à busca ao terminar.

INIT\_SP é a proposta de escrita de 227 em R29 depois da liberação do reset e antes da primeira busca. Essa solução ainda deve ser confrontada com o requisito de R29=227 enquanto o reset está ativo, conforme o Apêndice~\ref{ap:pendencias}. Na captura de \texttt{push}, A recebe o SP antigo e B o dado antigo de R[rs]; na de \texttt{pop}, A recebe o SP antigo. Os seletores precisam estar estáveis antes da borda que grava A/B. Se a classificação da instrução só terminar nessa borda, acrescentar um estado específico de captura.

\subsection{Sequência de \texttt{push} e \texttt{pop}}\label{sec:estados-pilha}

\begin{table}[H]\centering\footnotesize
\caption{Estados de \texttt{push} após a captura dos operandos.}\label{tab:push-estados}
\begin{tabularx}{\linewidth}{@{}p{2.5cm}p{4cm}X@{}}\toprule
Estado & Efeito na borda & Controles essenciais \\\midrule
PUSH\_DEC & $\mathrm{ALUOut}\gets A-4$ & ALUSrcA=A, ALUSrcB=4, F=010, ALUOutWrite=1. \\
PUSH\_SP & $R29\gets\mathrm{ALUOut}$ & RegDst=3, MemtoReg=0, RegWrite=1. \\
PUSH\_STORE & $M_{32}[\mathrm{ALUOut}]\gets B$ & IorD=1, DataSource=2, wr=1. \\\bottomrule
\end{tabularx}
\end{table}

A e B não são reescritos após a captura. Assim, \texttt{push R29} armazena o valor antigo do SP, preservado em B, apesar de R29 ter sido atualizado antes da escrita da memória. Durante esses estados, PCWrite, PCWriteCond e IRWrite permanecem desabilitados.

\begin{table}[H]\centering\small
\caption{Estados de \texttt{pop} após a captura de SP.}\label{tab:pop-estados}
\begin{tabularx}{\linewidth}{@{}p{2.5cm}p{4cm}X@{}}\toprule
Estado & Efeito na borda & Controles essenciais \\\midrule
POP\_ADDR & $\mathrm{ALUOut}\gets A$ & ALUSrcA=A, F=000, ALUOutWrite=1. \\
POP\_REQ & Requisitar leitura & IorD=1, wr=0, MDRWrite=0. \\
POP\_WAIT & Aguardar dado válido & Endereço estável, MDRWrite=0. \\
POP\_CAP & $\mathrm{MDR}\gets\mathrm{DataOut}$ & IorD=1, wr=0, MDRWrite=1. \\
POP\_RD & $R[rd]\gets\mathrm{MDR}$ & RegDst=1, MemtoReg=1, RegWrite=1. \\
POP\_INC & $\mathrm{ALUOut}\gets A+4$ & ALUSrcA=A, ALUSrcB=4, F=001, ALUOutWrite=1. \\
POP\_SP & $R29\gets\mathrm{ALUOut}$ & RegDst=3, MemtoReg=0, RegWrite=1. \\\bottomrule
\end{tabularx}
\end{table}

REQ, WAIT e CAP mantêm o mesmo endereço e representam uma leitura conservadora da memória síncrona. O \texttt{Memoria.vhd} registra o controle de rotação em dois estágios; a forma de onda da integração deve mostrar em qual borda a palavra desejada fica válida em \texttt{Dataout} e quantos estados são necessários antes de MDRWrite. A conserva SP antigo até POP\_INC. POP\_RD e POP\_SP usam bordas diferentes porque o banco dispõe de uma única porta de escrita.

\subsection{Exemplo de execução e bytes}\label{subsec:exemplo-pilha}

Considere R29=227 e R8=\texttt{0xAABBCCDD}. Em \texttt{push R8}, A conserva 227, B conserva a palavra de R8, a ALU produz 223, R29 recebe 223 e a memória recebe \texttt{0xAABBCCDD} em $M_{32}[223]$. Em \texttt{pop R9}, a memória lê a partir de 223, MDR conserva a palavra, R9 a recebe e R29 volta a 227. A memória não precisa ser apagada: o avanço de SP remove o elemento logicamente da pilha.

\begin{table}[H]\centering
\caption{Evolução esperada no exemplo da especificação.}\label{tab:exemplo-pilha}
\begin{tabular}{@{}lccc@{}}\toprule
Momento & R29 & R8 & R9 \\\midrule
Antes de \texttt{push R8} & 227 & \texttt{0xAABBCCDD} & --- \\
Após \texttt{push R8} & 223 & \texttt{0xAABBCCDD} & --- \\
Após \texttt{pop R9} & 227 & \texttt{0xAABBCCDD} & \texttt{0xAABBCCDD} \\\bottomrule
\end{tabular}
\end{table}

\subsection{Exceções e retorno}\label{subsec:excecoes}

Em opcode inexistente, overflow ou divisão por zero, a UC salva em EPC o endereço $P$ da instrução causadora. Depois lê o byte localizado em 253, 254 ou 255, estende esse byte com zeros para 32 bits e carrega o resultado em PC. Os números 253--255 são endereços que armazenam o vetor, não o destino direto do PC. A especificação menciona \texttt{rte} como retorno $PC\gets EPC$, mas não a codifica entre as 23 instruções; uma eventual implementação de \texttt{rte} deve ser decidida separadamente, sem alterar indevidamente o catálogo~\cite{especificacao}.

\begin{landscape}
\FiguraFigma{mapa-estados-completo.png}{Visão geral do mapa de estados e do contexto de controle exportados do FigJam.}{fig:mapa-figma}{0.98}
\FiguraFigma{mapa-estados-recorte.png}{Recorte ampliado do mapa de estados e das transições da UC.}{fig:mapa-recorte}{0.98}
\end{landscape}

\section{Conjunto de simulações}\label{sec:simulacoes}

Carlos Henrique Gonçalves dos Santos (CC, \texttt{chgs3}) foi responsável pela parte de simulação do projeto. Esta seção está preparada para incorporar as formas de onda e a análise da implementação ao texto. A especificação exige ao menos uma simulação de cada uma das 23 instruções, com descrição dos sinais, imagem Quartus e explicação de todos os passos. Em cada captura devem aparecer clock, reset, estado da UC, PC, EPC, MDR, IR, R29, R31 e os registradores envolvidos, além dos sinais que permitem compreender a instrução~\cite{especificacao}. Os campos \emph{resultado observado} permanecem abertos até a inserção das evidências correspondentes.

\subsection{Programas de apoio recebidos na parte 3}\label{subsec:programas-pt3}

A pasta \texttt{mips - alunos - 2026.2} fornece os programas da Tabela~\ref{tab:asm-pt3}. Eles são pontos de partida para gerar o conteúdo de \texttt{instrucoes.mif} e observar sequências reais da CPU; sua presença não equivale a resultados de simulação. O \texttt{LEIA-ME.txt} da parte 3 indica um conversor MIPS--MIF como apoio, mas cada imagem de memória deve ser conferida quanto às codificações específicas do projeto, sobretudo \texttt{push} e \texttt{pop}~\cite{docpt3}.

\begin{table}[H]\centering\small
\caption{Programas Assembly fornecidos e finalidade principal.}\label{tab:asm-pt3}
\begin{tabularx}{\linewidth}{@{}p{4cm}X@{}}\toprule
Arquivo & Situações que pode exercitar \\\midrule
\texttt{ula.asm}, \texttt{addi\_overflow.asm} & Aritmética, lógica, comparação e overflow de \texttt{addi}. \\
\texttt{branch.asm}, \texttt{jumps\_break.asm} & \texttt{beq}/\texttt{bne} tomados e não tomados; \texttt{j}, \texttt{jal} e \texttt{jr}. \\
\texttt{loads.asm}, \texttt{stores.asm} & \texttt{lw}/\texttt{lb} e \texttt{sw}/\texttt{sb} com dados em memória. \\
\texttt{shifts.asm} & \texttt{sll} e \texttt{sra}. \\
\texttt{mult.asm}, \texttt{div.asm} & Produto, quociente, resto, HI/LO e divisão por zero. \\
\texttt{push\_pop.asm} & Um \texttt{push} seguido de um \texttt{pop}; complementar com LIFO e \texttt{push R29}. \\
\texttt{excecoes.asm} & Cenários de divisão, overflow e opcode não reconhecido; validar a montagem do último estímulo. \\\bottomrule
\end{tabularx}
\end{table}

Os programas não substituem a exigência de uma captura e de uma análise para cada uma das 23 instruções. Também faltam, como estímulos explícitos, os casos de byte em todas as posições, palavra no endereço 223, pilha com duas entradas, limites de deslocamento e temporização de cada exceção. Esses cenários devem ser acrescentados pela equipe durante a validação.

\subsection{Convenção de preenchimento}\label{subsec:convencao-simulacao}

Para cada instrução: (1) registrar valores iniciais e programa de teste; (2) definir sinais visíveis e a escala temporal; (3) inserir uma captura legível; (4) percorrer os estados na ordem real, indicando o valor antes e depois de cada borda; (5) comparar valores medidos com os esperados; (6) declarar aprovação ou falha com causa conhecida. Os nomes das capturas podem ser inseridos diretamente no campo de figura de cada subseção ou associados a novos arquivos na pasta \texttt{figuras/}. Não usar o mapa de estados como substituto de uma forma de onda.

\subsection{Casos prioritários da pilha}\label{subsec:prioridade-pilha}

\begin{table}[H]\centering\small
\caption{Testes mínimos sugeridos para a extensão da pilha.}\label{tab:casos-pilha}
\begin{tabularx}{\linewidth}{@{}p{3.1cm}X X@{}}\toprule
Caso & Preparação & Resultado esperado \\\midrule
Inicialização & Liberar reset e entrar em INIT\_SP. & R29=227 antes da primeira busca. \\
\texttt{push R8} & R29=227; R8=\texttt{0xAABBCCDD}. & R29=223; $M_{32}[223]$ recebe a palavra. \\
\texttt{pop R9} & R29=223; $M_{32}[223]$ contém a palavra. & R9 recebe a palavra; R29=227. \\
LIFO & Dois \texttt{push}, seguidos de dois \texttt{pop}. & SP: 227, 223, 219, 223, 227; ordem invertida dos dados. \\
\texttt{push R29} & R29=227. & Memória em 223 recebe 227, não 223. \\
Leitura síncrona & Observar REQ, WAIT e CAP. & Endereço estável; MDRWrite apenas em CAP. \\
\emph{Little-endian} & Escrever \texttt{0x12345678} em 223. & Bytes 223--226: \texttt{78 56 34 12}. \\\bottomrule
\end{tabularx}
\end{table}

\subsection{Instruções do formato R}\label{subsec:sims-r}
\SimulacaoInstrucao{add}{add}{Verificar soma com sinal, escrita em rd e tratamento de overflow.}{rd recebe a soma válida; em overflow, a exceção segue a política da UC.}
\SimulacaoInstrucao{and}{and}{Verificar a operação lógica bit a bit.}{rd recebe a conjunção dos 32 pares de bits.}
\SimulacaoInstrucao{div}{div}{Verificar quociente, resto, conclusão e divisão por zero.}{LO recebe quociente, HI recebe resto; divisor zero sinaliza exceção.}
\SimulacaoInstrucao{mult}{mult}{Verificar Booth com sinal, produto de 64 bits e latência.}{HI/LO formam o produto; conclusão ocorre dentro do limite da especificação.}
\SimulacaoInstrucao{jr}{jr}{Verificar salto por registrador.}{PC recebe o conteúdo de R[rs].}
\SimulacaoInstrucao{mfhi}{mfhi}{Verificar leitura de HI sem alterá-lo.}{rd recebe o valor armazenado em HI.}
\SimulacaoInstrucao{mflo}{mflo}{Verificar leitura de LO sem alterá-lo.}{rd recebe o valor armazenado em LO.}
\SimulacaoInstrucao{sll}{sll}{Verificar deslocamento lógico, inclusive quantidades 0 e 31.}{rd recebe R[rt] deslocado à esquerda, com zeros inseridos.}
\SimulacaoInstrucao{slt}{slt}{Verificar comparação com sinal e extremos.}{rd recebe exatamente 0 ou 1.}
\SimulacaoInstrucao{sra}{sra}{Verificar deslocamento aritmético e preservação do sinal.}{rd recebe R[rt] deslocado à direita com replicação do bit de sinal.}
\SimulacaoInstrucao{sub}{sub}{Verificar subtração com sinal e overflow.}{rd recebe a diferença válida; overflow desvia ao tratamento.}
\SimulacaoInstrucao{push}{push}{Verificar captura de R[rs] antigo, decremento de SP e escrita da palavra.}{SP diminui quatro; a memória recebe o valor preservado em B.}
\SimulacaoInstrucao{pop}{pop}{Verificar leitura síncrona, escrita de rd e incremento posterior de SP.}{rd recebe a palavra do topo antigo; R29 aumenta quatro.}

\subsection{Instruções do formato I}\label{subsec:sims-i}
\SimulacaoInstrucao{addi}{addi}{Verificar extensão de sinal do imediato e soma.}{rt recebe a soma válida; overflow é tratado.}
\SimulacaoInstrucao{beq}{beq}{Verificar condições verdadeira/falsa e offset positivo/negativo.}{PC recebe o alvo relativo a P+4 somente quando R[rs]=R[rt].}
\SimulacaoInstrucao{bne}{bne}{Verificar a inversão da condição de igualdade.}{PC recebe o alvo somente quando R[rs] difere de R[rt].}
\SimulacaoInstrucao{lb}{lb}{Verificar as quatro posições do byte e extensão de sinal.}{rt recebe o byte endereçado estendido para 32 bits.}
\SimulacaoInstrucao{lui}{lui}{Verificar seleção do imediato e deslocamento de 16 posições.}{rt recebe imediato nos 16 bits altos e zeros nos baixos.}
\SimulacaoInstrucao{lw}{lw}{Verificar endereço efetivo, latência da memória e MDR.}{rt recebe a palavra lida após a captura válida.}
\SimulacaoInstrucao{sb}{sb}{Verificar escrita de byte sem perda dos vizinhos.}{Apenas o byte endereçado muda na palavra da memória.}
\SimulacaoInstrucao{sw}{sw}{Verificar que o dado gravado vem de R[rt] via B.}{A palavra em EA recebe o valor de R[rt].}

\subsection{Instruções do formato J}\label{subsec:sims-j}
\SimulacaoInstrucao{j}{j}{Verificar formação do alvo e atualização do PC.}{PC recebe o destino formado pelos bits altos e target deslocado.}
\SimulacaoInstrucao{jal}{jal}{Verificar alvo e endereço de retorno após definir a convenção.}{PC recebe o alvo; R31 recebe o link acordado para esta CPU.}

\subsection{Simulações isoladas das entidades implementadas}\label{subsec:sims-entidades}

A especificação dispensa simulação isolada dos componentes fornecidos e dos multiplexadores, mas exige imagem e explicação para cada unidade construída pela equipe. Para cada entidade implementada, duplicar o modelo abaixo e preencher nome, portas, estímulos, resultado esperado, resultado observado e figura.

\subsubsection{Unidade de controle}\label{sim-ent:uc}
\noindent\textbf{Entradas e saídas:} \textit{Preencher de acordo com o HDL.}\par
\noindent\textbf{Cenários:} decodificação das 23 instruções, transições de retorno à busca, reset, exceções e estados da pilha.\par
\noindent\textbf{Forma de onda e análise:} \textit{Inserir captura e explicar cada transição relevante.}\par

\subsubsection{Multiplicador}\label{sim-ent:mult}
\noindent\textbf{Portas e algoritmo:} \textit{Preencher conforme a unidade implementada.}\par
\noindent\textbf{Cenários:} sinais dos operandos, produto de 64 bits, início, conclusão, 32--40 ciclos e reset interno.\par
\noindent\textbf{Forma de onda e análise:} \textit{Inserir captura e conferir HI/LO.}\par

\subsubsection{Divisor}\label{sim-ent:div}
\noindent\textbf{Portas e algoritmo:} \textit{Preencher conforme a unidade implementada.}\par
\noindent\textbf{Cenários:} quociente, resto, sinais, divisor zero, conclusão e reset interno.\par
\noindent\textbf{Forma de onda e análise:} \textit{Inserir captura e conferir HI/LO.}\par

\subsubsection{Outras entidades próprias}\label{sim-ent:outras}
\noindent\textbf{Entidades a incluir:} \textit{Listar apenas as efetivamente implementadas pela equipe, por exemplo MergeBytes, comparador ou lógica adicional de exceções. Criar uma subsubseção e figura para cada uma.}\par

\section{Conclusão}\label{sec:conclusao}

O projeto documental organiza uma CPU multiciclo para as 23 instruções da especificação 2026.2 e explicita como a extensão de pilha reutiliza o banco, a ALU e a memória. \texttt{push} conserva o operando antes de atualizar R29; \texttt{pop} lê a palavra no SP antigo antes de incrementar o ponteiro. A sequência proposta reproduz o exemplo 227 $\to$ 223 $\to$ 227 e reserva ciclos separados para a leitura síncrona e as duas escritas do \texttt{pop}.

A parte 3 disponibiliza componentes VHDL, inicialização da memória e programas Assembly que tornam mais precisa a análise das interfaces e do plano de testes. As imagens do circuito e do mapa de estados já foram incorporadas. A avaliação do funcionamento efetivo no relatório ainda depende de documentar a integração da CPU e da UC e inserir as capturas e análises das 23 instruções e das entidades construídas pela equipe. As subseções da Seção~\ref{sec:simulacoes} foram preparadas para receber esses dados sem confundir expectativa com medição. Antes da entrega, a equipe deve resolver as pendências do Apêndice~\ref{ap:pendencias} e conferir as figuras contra a implementação final.

\begin{thebibliography}{9}
\bibitem{especificacao} UNIVERSIDADE FEDERAL DE PERNAMBUCO. \emph{Infra-Estrutura de Hardware: Especificação de Projeto 2026.2}. Recife, 2026. Arquivo local: \texttt{Projeto/Especificacao Projeto 2026-2.pdf}.
\bibitem{docpt1} EQUIPE EQ01. \emph{Documentação PT.1: análise, catálogo de instruções, arquitetura, estados, decisões e validação}. Pasta \texttt{Projeto/Documentacao PT.1}, 2026.
\bibitem{docpt2} EQUIPE EQ01. \emph{Documentação PT.2: roteiro e guia lógico do mapa de estados}. Pasta \texttt{Projeto/Documentacao PT.2}, 2026.
\bibitem{docpt3} IF674. \emph{Documentação PT.3: componentes VHDL, arquivo inicial da memória e programas Assembly de teste}. Pasta \texttt{Projeto/Documentação PT.3}, 2026.
\bibitem{aulas} IF674. \emph{Slides consolidados de instruções, processador monociclo, processador multiciclo e desempenho}. Pasta \texttt{Aulas/Slides consolidados}, 2026.
\bibitem{figma} EQUIPE EQ01. \emph{Projeto IH 2026.2 — quadro do circuito e mapa de estados}. FigJam. \url{https://www.figma.com/board/hILCPGWNEfBjOQvYNMzxR0/ProjetoIH-2K26.2--EQ01?node-id=0-1}. Acesso em outubro de 2026.
\end{thebibliography}

\appendix
\section{Catálogo das 23 instruções}\label{ap:catalogo}

As codificações seguem as Tabelas 1--3 da especificação 2026.2. Em todas as instruções R, \texttt{opcode=0x00}. \texttt{push} e \texttt{pop} são extensões desta disciplina e podem não ser aceitas por montadores MIPS genéricos.

\small
\begin{longtable}{@{}p{1.42cm}p{3.45cm}p{2.25cm}p{6.93cm}@{}}
\caption{Instruções exigidas e efeito arquitetural principal.}\label{tab:catalogo}\\
\toprule Formato & Instrução & Código & Efeito principal \\\midrule
\endfirsthead
\toprule Formato & Instrução & Código & Efeito principal \\\midrule
\endhead
\bottomrule
\endfoot
R & \texttt{add rd,rs,rt} & funct 0x20 & Soma com sinal, destino rd e controle de overflow. \\
R & \texttt{and rd,rs,rt} & funct 0x24 & AND bit a bit em rd. \\
R & \texttt{div rs,rt} & funct 0x1A & Quociente em LO, resto em HI, erro de divisor zero. \\
R & \texttt{mult rs,rt} & funct 0x18 & Produto com sinal de 64 bits em HI/LO. \\
R & \texttt{jr rs} & funct 0x08 & PC recebe R[rs]. \\
R & \texttt{mfhi rd} & funct 0x10 & rd recebe HI. \\
R & \texttt{mflo rd} & funct 0x12 & rd recebe LO. \\
R & \texttt{sll rd,rt,shamt} & funct 0x00 & Deslocamento lógico à esquerda. \\
R & \texttt{slt rd,rs,rt} & funct 0x2A & Comparação com sinal; rd recebe 0 ou 1. \\
R & \texttt{sra rd,rt,shamt} & funct 0x03 & Deslocamento aritmético à direita. \\
R & \texttt{sub rd,rs,rt} & funct 0x22 & Subtração com sinal e controle de overflow. \\
R & \texttt{push rs} & funct 0x05 & R29 diminui quatro; R[rs] antigo é gravado no novo topo. \\
R & \texttt{pop rd} & funct 0x06 & rd recebe a palavra do topo antigo; R29 aumenta quatro. \\
I & \texttt{addi rt,rs,imm} & opcode 0x08 & Soma com imediato estendido com sinal. \\
I & \texttt{beq rs,rt,off} & opcode 0x04 & Desvio relativo se R[rs]=R[rt]. \\
I & \texttt{bne rs,rt,off} & opcode 0x05 & Desvio relativo se R[rs]$\ne$R[rt]. \\
I & \texttt{lb rt,off(rs)} & opcode 0x20 & Leitura de byte e extensão de sinal. \\
I & \texttt{lui rt,imm} & opcode 0x0F & Imediato nos 16 bits altos. \\
I & \texttt{lw rt,off(rs)} & opcode 0x23 & Leitura de palavra em endereço efetivo. \\
I & \texttt{sb rt,off(rs)} & opcode 0x28 & Escrita de byte com preservação dos vizinhos. \\
I & \texttt{sw rt,off(rs)} & opcode 0x2B & Escrita da palavra de R[rt] em endereço efetivo. \\
J & \texttt{j target} & opcode 0x02 & Salto incondicional para alvo formado. \\
J & \texttt{jal target} & opcode 0x03 & Salto com gravação de endereço de retorno em R31. \\
\end{longtable}
\normalsize

\section{Decisões a confirmar antes da entrega}\label{ap:pendencias}

Os itens abaixo separam requisitos explícitos, convenções provisórias e verificações que só podem ser encerradas com o HDL ou com orientação da disciplina. Sua presença evita que o relatório atribua certeza a comportamentos ainda não demonstrados.

\begin{longtable}{@{}p{1.1cm}p{3.1cm}p{10.3cm}@{}}
\caption{Pendências de integração e validação.}\label{tab:pendencias}\\
\toprule ID & Tema & Conferência necessária \\\midrule
\endfirsthead
\toprule ID & Tema & Conferência necessária \\\midrule
\endhead
\bottomrule
\endfoot
D01 & \texttt{jal} & Confirmar link e existência de \emph{delay slot}. A documentação assume provisoriamente $P+4$, sem \emph{delay slot}. \\
D02 & Reset de R29 & Conciliar exigência de SP=227 durante reset com o banco fornecido, cujo reset assíncrono limpa os registradores. INIT\_SP só escreve após a liberação. \\
D03 & Endereço 223 & A rotação dos bancos de \texttt{Memoria.vhd} prevê endereço desalinhado; demonstrar em simulação que a palavra em 223 usa exatamente os bytes 223--226. \\
D04 & \texttt{pop R29} & Definir a prioridade entre a escrita do dado lido em rd=29 e a atualização final de SP. A sequência proposta faz o incremento prevalecer. \\
D05 & Divisão extrema & Definir comportamento de \texttt{0x80000000}/$-1$, cujo quociente não cabe em 32 bits com sinal. \\
D06 & \texttt{rte} & A especificação menciona retorno via EPC, mas não fornece sua codificação entre as 23 instruções. \\
D07 & \texttt{lb}/\texttt{slt} & Confirmar extensão de sinal de byte e formação de 0/1 em 32 bits para comparação com sinal. \\
D08 & Latência de memória & Confrontar REQ/WAIT/CAP com a temporização real da memória e de MDR antes de otimizar estados. \\
D09 & Outras instruções & Conferir origem B de \texttt{sw}, condição invertida de \texttt{bne}, bytes de \texttt{sb}, alvo de saltos, conclusão de MULT/DIV e vetores de exceção. \\
D10 & Integração da parte 3 & Integrar os módulos fornecidos sem modificar seu código, completar CPU, UC, MULT/DIV e lógica de bytes; conferir portas e sinais reais. \\
D11 & Programas de teste & Gerar o MIF dos programas Assembly e complementar os estímulos para cobrir as 23 instruções, LIFO, endereço 223 e exceções. \\
\end{longtable}

\end{document}
